\documentclass[10pt,a4paper]{amsart}
\usepackage[margin=19mm]{geometry}
\usepackage{amsmath,amssymb,mathtools}
\usepackage[scr=rsfs,cal=euler]{mathalfa}
\usepackage{xfrac}
\usepackage[unicode,hidelinks,bookmarks=false]{hyperref}
\newcommand{\ie}{i.e.\ }
\newcommand{\cf}{cf.\ }
\newcommand{\resp}{respectively\ }
\newcommand{\Subsection}{\subsection}
\providecommand{\nobreakdash}{\nobreak\hbox{-}}
\setlength{\emergencystretch}{2em}
\multlinegap0pt
\usepackage{amssymb}
\usepackage{bm}
\newtheorem{prop}{Proposition}
\newcommand{\dd}{\,\mathrm{d}}
\newcommand{\la}{\langle}
\newcommand{\pr}{\partial}
\newcommand{\ra}{\rangle}
\title{Inverse boundary value problems for certain doubly nonlinear parabolic and elliptic equations}
\author{C\u{a}t\u{a}lin I. C\^arstea, Tuhin Ghosh}
\date{Source: arXiv:2603.08297v1 (CC BY 4.0)}
\begin{document}
\maketitle

\noindent\emph{Source and licence.} Inverse boundary value problems for certain doubly nonlinear parabolic and elliptic equations. Version arXiv:2603.08297v1 (CC BY 4.0), distributed by arXiv under the Creative Commons Attribution 4.0 International licence, http://creativecommons.org/licenses/by/4.0/, as recorded in the arXiv licence metadata for this exact version and shown on the abstract page https://arxiv.org/abs/2603.08297v1. Full licence notice: \texttt{licenses/arxiv-2603.08297-CC-BY-4.0.txt}.\par\smallskip

\noindent\textit{Editorial scope.} Counted unit: the article's prop prop:small-data (source0.tex lines 626--642). The article's complete original proof of it is delivered with it (source0.tex lines 644--722, one \texttt{proof} environment of 79 lines). No mathematics is written, repaired or re-derived: the statement and the proof are the article's own lines, in the article's own notation, and no retained line is substituted or re-flowed.  Retained uncounted as the source-local context the proof needs: lines 500--507; lines 523--551; lines 584--587; lines 596--598; lines 613--620; lines 622--624.  Nothing was omitted from any retained span, so the package carries no omission record.  No retained line contains a citation, so the wrapper carries no bibliography and no bibliography entry.  No retained line contains a figure environment, an image-inclusion command or drawing code, so no image is delivered and the package declares no asset dependency.  Editorial formatting only, in addition: an AMS \texttt{amsart} preamble that restates the article's own declarations and macros used by the retained lines, namely the environments \texttt{prop} on separate counters, together with the standard packages those lines need.  The retained environments print in this document as Proposition 1 in order of appearance, whereas the article prints the same items with the numbering of its own full document.  The complete native source of the article as distributed in the arXiv e-print for this exact version is bundled unchanged as \texttt{sources/195994/source0.tex}.
\begin{equation}\label{eq:elliptic}
\left\{
\begin{aligned}
-\nabla\cdot\bigl(\gamma|\nabla w|^{p-2}\nabla w\bigr)+Vw^m&=0 &&\text{in }\Omega,\\
w&=g &&\text{on }\pr\Omega,
\end{aligned}
\right.
\end{equation}

\begin{prop}\label{prop:elliptic-forward}
Let $g\in W^{1,p}(\Omega)\cap L^{m+1}(\Omega)$ be nonnegative. Then \eqref{eq:elliptic} admits a unique weak solution
\[
w\in g+W^{1,p}_0(\Omega),\qquad w\ge0.
\]
Moreover, the following assertions hold.
\begin{enumerate}
\item The solution $w$ is the unique minimizer of the energy functional
\begin{equation}\label{eq:energy}
\mathcal E[v]
=\int_\Omega\left(\frac{1}{p}\gamma|\nabla v|^p+\frac{1}{m+1}Vv^{m+1}\right)\,dx
\end{equation}
over the affine space $g+W^{1,p}_0(\Omega)$.
\item There exists a constant $C$, depending only on $\Omega$, $p$, $m$, and upper and lower bounds for $\gamma$ and $V$, such that
\begin{equation}\label{eq:elliptic-energy-bound}
\|w\|_{W^{1,p}(\Omega)}+\|w\|_{L^{m+1}(\Omega)}
\le C\bigl(1+\|g\|_{W^{1,p}(\Omega)}+\|g\|_{L^{m+1}(\Omega)}\bigr).
\end{equation}
\item If $g\in L^\infty(\pr\Omega)$, then
\begin{equation}\label{eq:elliptic-max-principle}
\|w\|_{L^\infty(\Omega)}\le \|g\|_{L^\infty(\pr\Omega)}.
\end{equation}
\item If $g\in C^{1,\beta}(\pr\Omega)$ for some $\beta\in(0,1)$, then $w\in C^{1,\beta}(\overline\Omega)$ and
\begin{equation}\label{eq:lieberman-estimate}
\|w\|_{C^{1,\beta}(\overline\Omega)}
\le C\bigl(1+\|g\|_{C^{1,\beta}(\pr\Omega)}\bigr).
\end{equation}
\end{enumerate}
\end{prop}

\begin{equation}\label{eq:p-laplace}
\nabla\cdot\bigl(\gamma|\nabla v|^{p-2}\nabla v\bigr)=0
\qquad \text{in }\Omega,
\end{equation}

\begin{equation}\label{eq:Taylor-J}
J_j(\zeta)=J_j(\xi)+\sum_{k=1}^n(\zeta_k-\xi_k)\int_0^1 \partial_{\xi_k}J_j\bigl(\xi+t(\zeta-\xi)\bigr)\,\dd t,
\end{equation}

\begin{equation}\label{eq:w-lambda-small}
\left\{
\begin{aligned}
-\nabla\cdot\bigl(\gamma|\nabla w_\lambda|^{p-2}\nabla w_\lambda\bigr)+Vw_\lambda^m&=0 &&\text{in }\Omega,\\
 w_\lambda&=\lambda v &&\text{on }\pr\Omega.
\end{aligned}
\right.
\end{equation}

\begin{equation}\label{eq:ansatz-small}
w_\lambda=\lambda\bigl(v+\lambda^{m-p+1}R_\lambda\bigr).
\end{equation}

\begin{prop}\label{prop:small-data}
Assume that $m>p-1$. Then the family $R_\lambda$ is bounded in $C^{1,\beta}(\overline\Omega)$, and there exists $R\in C^{1,\beta}(\overline\Omega)$ such that $R_\lambda\to R$ in $C^1(\overline\Omega)$ as $\lambda\to0^+$. Moreover, $R$ solves
\begin{equation}\label{eq:R-equation}
\left\{
\begin{aligned}
\nabla\cdot\bigl(A[v]\nabla R\bigr)&=Vv^m &&\text{in }\Omega,\\
R&=0 &&\text{on }\pr\Omega,
\end{aligned}
\right.
\end{equation}
and for every $\omega\in C^\infty(\overline\Omega)$,
\begin{multline}\label{eq:expansion-small}
\la\omega,\Lambda_{\gamma,V}(\lambda v|_{\pr\Omega})\ra
=\lambda^{p-1}\int_\Omega \gamma|\nabla v|^{p-2}\nabla v\cdot\nabla\omega\,\dd x\\
+\lambda^m\int_\Omega\Bigl(\nabla\omega\cdot A[v]\nabla R+V\omega v^m\Bigr)\,\dd x+o(\lambda^m).
\end{multline}
\end{prop}

\begin{proof}
By Proposition~\ref{prop:elliptic-forward}, the family
\[
v+\lambda^{m-p+1}R_\lambda=\lambda^{-1}w_\lambda
\]
is bounded in $C^{1,\beta}(\overline\Omega)$ uniformly in $\lambda\in(0,1)$. By the Arzel\`a--Ascoli theorem, after passing to a subsequence there exists $\tilde v\in C^{1,\beta}(\overline\Omega)$ such that $\tilde v|_{\pr\Omega}=v|_{\pr\Omega}$ and
\[
v+\lambda^{m-p+1}R_\lambda\to \tilde v
\qquad\text{in }C^1(\overline\Omega)
\]
as $\lambda\to0^+$.

We next identify the limit. Since $w_\lambda$ solves \eqref{eq:w-lambda-small}, we have
\[
\lambda^{1-p}\Bigl(-\nabla\cdot\bigl(\gamma|\nabla w_\lambda|^{p-2}\nabla w_\lambda\bigr)+Vw_\lambda^m\Bigr)=0
\]
in the sense of distributions. Using \eqref{eq:ansatz-small}, the convergence above, and the fact that $m>p-1$, we obtain
\[
\mathscr D'(\Omega)\text{-}\lim_{\lambda\to0^+}\lambda^{1-p}Vw_\lambda^m=0
\]
and
\[
\mathscr D'(\Omega)\text{-}\lim_{\lambda\to0^+}\lambda^{1-p}\nabla\cdot\bigl(\gamma|\nabla w_\lambda|^{p-2}\nabla w_\lambda\bigr)
=\nabla\cdot\bigl(\gamma|\nabla\tilde v|^{p-2}\nabla\tilde v\bigr).
\]
Therefore $\tilde v$ satisfies the same weighted $p$-Laplace equation \eqref{eq:p-laplace} as $v$, with the same Dirichlet data. By uniqueness for the weighted $p$-Laplace equation with boundary value $v|_{\pr\Omega}$, we conclude that $\tilde v=v$. Hence
\[
\lambda^{m-p+1}R_\lambda\to0
\qquad\text{in }C^1(\overline\Omega)
\]
as $\lambda\to0^+$. Since the limit is independent of the subsequence chosen, the convergence holds for the full family.

Applying Taylor's formula \eqref{eq:Taylor-J} with
\[
\xi=\nabla v,
\qquad
\zeta=\nabla v+\lambda^{m-p+1}\nabla R_\lambda,
\]
we obtain
\[
\gamma|\nabla w_\lambda|^{p-2}\nabla w_\lambda
=\lambda^{p-1}\gamma|\nabla v|^{p-2}\nabla v+\lambda^m A_\lambda\nabla R_\lambda,
\]
where
\[
(A_\lambda)_{jk}=\int_0^1\gamma\,\pr_{\xi_k}J_j\bigl(\nabla v+t\lambda^{m-p+1}\nabla R_\lambda\bigr)\,\dd t.
\]
Since $v$ solves \eqref{eq:p-laplace}, it follows that
\[
\nabla\cdot(A_\lambda\nabla R_\lambda)=V\bigl(v+\lambda^{m-p+1}R_\lambda\bigr)^m
\qquad\text{in }\Omega,
\]
with $R_\lambda|_{\pr\Omega}=0$. In particular, because $\lambda^{m-p+1}R_\lambda\to0$ in $C^1(\overline\Omega)$ and $\nabla v$ has no zeros in $\overline\Omega$, the coefficient matrices $A_\lambda$ are uniformly elliptic for all sufficiently small $\lambda$, with ellipticity constants independent of $\lambda$. Moreover,
\[
V\bigl(v+\lambda^{m-p+1}R_\lambda\bigr)^m\to Vv^m
\qquad\text{in }C(\overline\Omega).
\]
Schauder estimates therefore imply that $R_\lambda$ is bounded in $C^{1,\beta}(\overline\Omega)$ uniformly in $\lambda$. Passing again to a subsequence, we obtain $R\in C^{1,\beta}(\overline\Omega)$ such that $R_\lambda\to R$ in $C^1(\overline\Omega)$, and the limit solves \eqref{eq:R-equation}. By uniqueness for \eqref{eq:R-equation}, the whole family converges to $R$.

Finally, by the weak definition of the Dirichlet-to-Neumann map,
\[
\la\omega,\Lambda_{\gamma,V}(\lambda v|_{\pr\Omega})\ra
=\int_\Omega \gamma|\nabla w_\lambda|^{p-2}\nabla w_\lambda\cdot\nabla\omega\,\dd x
+\int_\Omega V\omega w_\lambda^m\,\dd x.
\]
Using the expansion above for the flux term, together with the convergences $A_\lambda\to A[v]$ and $R_\lambda\to R$ in $C(\overline\Omega)$, we obtain
\begin{multline*}
\int_\Omega \gamma|\nabla w_\lambda|^{p-2}\nabla w_\lambda\cdot\nabla\omega\,\dd x
=\lambda^{p-1}\int_\Omega \gamma|\nabla v|^{p-2}\nabla v\cdot\nabla\omega\,\dd x\\[5pt]
+\lambda^m\int_\Omega \nabla\omega\cdot A[v]\nabla R\,\dd x+o(\lambda^m).
\end{multline*}
Also,
\[
w_\lambda^m=\lambda^m\bigl(v+\lambda^{m-p+1}R_\lambda\bigr)^m
=\lambda^m v^m+o(\lambda^m)
\qquad\text{uniformly in }\overline\Omega.
\]
Combining these two expansions yields \eqref{eq:expansion-small}.
\end{proof}

\end{document}
